\documentclass[main.tex]{subfiles}


\begin{document}

% \AddToShipoutPictureFG{%
% \begin{tikzpicture}[overlay,remember picture]
% \path (current page.south west) -- (current page.north east)
% node[midway,scale=1,color=gray,sloped,opacity=0.4] {\textnormal{Кравченко Игорь Игоревич\hspace{1cm} +79010144910\hspace{1cm} super.antig@yandex.ru}};
% \end{tikzpicture}
% }


% %from pagenumber to up
% \phantomsection 
% \hypertarget{toc}{}

 

 

 
   

\section{Сила трения}

\textbf{Трение}~--- это взаимодействие соприкасающихся тел, препятствующее их относительному движению. 

Возникновение трения можно объяснить следующим образом. Соприкасающиеся поверхности тел являются шероховатыми, они усеяны микроскопическими, незаметными глазу бугорками\footnote{Для наглядности на рисунке далее бугорки
сильно преувеличены.}
разных форм и размеров. Эти бугорки зацепляются друг за друга и мешают телам двигаться друг относительно друга.
\looseness=-1

\sbox{0}{%
    \begin{tikzpicture}[font=\small,            line cap=round,                             >={Stealth[inset=0pt,round,length=3mm, angle'=20]},vec/.style={>={Stealth[inset=0pt,round,length=3.5mm, angle'=20]}}]

\begin{scope}[rotate=30]

\path[spath/save=rpath,name path=one] decorate[decoration={random steps,segment length=3pt,amplitude=1pt}] {(0,0) -- (2,0) coordinate (le) -- (3,0) coordinate (ri) -- (5,0)};

%solid
\path[lightgray,draw=black,name path=two] (le) -- (2,.5) -- (3,.5) -- (ri);
\begin{scope}
    \clip[] [spath/use=rpath] --++ (150:{5*cos(30)}) -- (0,0);
    \filldraw[lightgray,draw=black] (2,-.15) -- (2,.5) -- (3,.5) -- (3,-.15) -- cycle;
    
\end{scope}

\filldraw[lightgray,semitransparent] [spath/use=rpath] --++ (0,-.3) --++ (-5,0) -- cycle;
\draw[] [spath/use=rpath];

%vec
\draw[->,red,thick,vec] (2.5,.25) --++ (-120:2) node[black,right] {$\vec F_\text{т}$};
\node at (2.5,.25) [circle,fill=red,inner sep=0pt,minimum size=1mm,label=above:] {};

\draw[->,red,thick,vec] (2.5,.25) --++ (0,{2*cos(30)}) node[black,left] {$\vec N$};

\draw[->,red,thick,vec] (2.5,.25) --++ ({2*cos(60)},0) node[black,above] {$\vec F_\text{тр}$};


\end{scope}


    \end{tikzpicture}%
}

{%
\setlength\intextsep{0pt}
\begin{wrapfigure}[12]{r}{\wd0}
    \centering
    \usebox{0}%
    \caption{Тело на наклонной плоскости}
    \label{fig:p41}
\end{wrapfigure}


Пусть тело покоится на шероховатой \emph{наклонной плоскости} (рис.\,\ref{fig:p41}). 

% \begin{figure}[H]
%     \centering

%     \begin{tikzpicture}[font=\small,            line cap=round,                             >={Stealth[inset=0pt,round,length=3mm, angle'=20]},vec/.style={>={Stealth[inset=0pt,round,length=3.5mm, angle'=20]}}]

% \begin{scope}[rotate=30]

% \path[spath/save=rpath,name path=one] decorate[decoration={random steps,segment length=3pt,amplitude=1pt}] {(0,0) -- (2,0) coordinate (le) -- (3,0) coordinate (ri) -- (5,0)};

% %solid
% \path[lightgray,draw=black,name path=two] (le) -- (2,.5) -- (3,.5) -- (ri);
% \begin{scope}
%     \clip[] [spath/use=rpath] --++ (150:{5*cos(30)}) -- (0,0);
%     \filldraw[lightgray,draw=black] (2,-.15) -- (2,.5) -- (3,.5) -- (3,-.15) -- cycle;
    
% \end{scope}

% \filldraw[lightgray,semitransparent] [spath/use=rpath] --++ (0,-.3) --++ (-5,0) -- cycle;
% \draw[] [spath/use=rpath];

% %vec
% \draw[->,red,thick,vec] (2.5,.25) --++ (-120:2) node[black,right] {$\vec F_\text{т}$};
% \node at (2.5,.25) [circle,fill=red,inner sep=0pt,minimum size=1mm,label=above:] {};

% \draw[->,red,thick,vec] (2.5,.25) --++ (0,{2*cos(30)}) node[black,left] {$\vec N$};

% \draw[->,red,thick,vec] (2.5,.25) --++ ({2*cos(60)},0) node[black,above] {$\vec F_\text{тр}$};


% \end{scope}


%     \end{tikzpicture}
% \caption{Тело на наклонной плоскости}
% \label{fig:p41}
% \end{figure}


Во"=первых, тело притягивается вниз к планете силой тяжести~$\vec F_\text{т}$; во"=вторых, на тело со стороны опоры действует сила~$\vec N$, перпендикулярная плоскости (независимо от того, шероховатая поверхность или гладкая). Из рис.\,\ref{fig:p41} видно, что в данном случае сумма сил $\vec F_\text{т} + \vec N$ не может быть равна нулю; и под воздействием только этих сил тело получило бы возможность соскальзывать вниз по наклонной плоскости. Однако шероховатость опоры обусловливает еще одну силу~--- \textbf{силу трения}~$\vec F_\text{тр}$, направленную против возможного движения тела (тело <<пытается>> начать скользить, но трение удерживает его на месте). Сила трения всегда направлена вдоль
поверхностей соприкасающихся тел и служит проявлением \emph{электромагнитного взаимодействия} тел.

}
% . ландсберг с 130
% баканина силы трения 1973 г. №9 квант 


Трение между телами, происходящее без жидкой или газообразной прослойки между ними, называется \textbf{сухим трением}. Принято выделять три основных вида сил такого трения\footnote{Эти силы, будучи приложенными к телу со стороны шероховатой поверхности, направлены против движения (или возможного движения) тела относительно этой поверхности.}.
\begin{enumerate}
    \item \textbf{Сила трения покоя} препятствует возникновению относительного движения тел:
    \begin{equation}
        F_\text{тр.\,п}\leqslant\mu N,
    \end{equation}
    где~$\mu$~--- \emph{коэффициент трения}, характеризующий шероховатость тел. 
    
    \emph{Пример:} тело покоится на шероховатой наклонной плоскости.

    \item \textbf{Сила трения скольжения} препятствует проскальзыванию поверхностей тел:
    \begin{equation}
        F_\text{тр.\,с}=\mu N.
    \end{equation}
    
    \emph{Пример:} тело скользит по шероховатой поверхности.

    \item \textbf{Сила трения качения} препятствует перекатыванию одного тела по другому.

    \emph{Пример:} тело катится по нетвердой опоре.
\end{enumerate}

Можно условно считать, что между телом, движущимся в жидкости или газе, и средой, окружающей тело, также происходит трение, которое называют \emph{сопротивлением среды}. Соответственно, против скорости тела направлена и \emph{сила сопротивления среды}, действующая на это тело.















 
\end{document}